%% ma: L10-B14-0079
%% lop: 10
%% chuong: 5
%% bai: 14
%% dang: 10.14.7
%% loai: TN4
%% muc: TH
%% dap_an: "A"
%% nguon: "chuyen-de-cac-so-dac-trung-cua-mau-so-lieu-khong-ghep-nhom-toan-10-chuong-trinh-moi.pdf (GV. Phan Nhật Linh) (P05), Bài 03 Dạng 2 Phần I Câu 26"
%% kien_thuc: "Lớp 10 Bài 13 (số trung bình), Bài 14 (phương sai, độ lệch chuẩn)"
%% trang_thai: chinh-thuc
%% ghi_chu: "Tự duyệt theo yêu cầu giáo viên 22/10/2026."
\begin{ex}
Nhiệt độ cao nhất của Hà Nội trong $7$ ngày liên tiếp của tháng 8 được ghi lại là: $34$; $34$; $36$; $35$; $33$; $31$; $30$ (đơn vị: $^\circ$C). Độ lệch chuẩn của mẫu số liệu thuộc khoảng nào?
\choice{\True $(1;2)$}{$(3;4)$}{$\left[2;\dfrac72\right]$}{$\left(0;\dfrac34\right)$}
\loigiai{$\overline x=\dfrac{233}7\approx33{,}29$; $s^2\approx3{,}92$; $s\approx1{,}98\in(1;2)$. Chọn A.}
\end{ex}
